\begin{statementbox}
\textbf{\textcolor{CStatement}{F014｜分段函数怎样才算整体递增或递减}}
\begin{center}\begin{tikzpicture}[font=\small,>=stealth,line join=round]
\draw[gray!40,rounded corners=3pt] (0,0) rectangle (14.2,7.4);
\node at (7.1,6.800000000000001) {两段斜率都为正，交接方向决定整体是否递增};
\begin{scope}[shift={(2.1,3.1)},x=0.65cm,y=0.65cm]\draw[->,gray] (-2,0)--(2.3,0) node[right]{$x$};\draw[->,gray] (0,-2)--(0,4.3) node[above]{$y$};\begin{scope}\clip (-2,-2) rectangle (2,4);\draw[blue,thick,domain=-2:0,samples=80] plot(\x,{\x});\draw[orange!85!black,thick,domain=0.0001:2,samples=100] plot(\x,{\x+2});\end{scope}\draw[draw=blue,fill=blue] (0,0) circle(2pt);\draw[draw=orange!85!black,fill=white] (0,2) circle(2pt);\end{scope}\node at (2.1,1.35) {向上跳：严格增};\begin{scope}[shift={(6.699999999999999,3.1)},x=0.65cm,y=0.65cm]\draw[->,gray] (-2,0)--(2.3,0) node[right]{$x$};\draw[->,gray] (0,-2)--(0,4.3) node[above]{$y$};\begin{scope}\clip (-2,-2) rectangle (2,4);\draw[blue,thick,domain=-2:0,samples=80] plot(\x,{\x});\draw[orange!85!black,thick,domain=0.0001:2,samples=100] plot(\x,{\x});\end{scope}\draw[draw=blue,fill=blue] (0,0) circle(2pt);\draw[draw=blue,fill=blue] (0,0) circle(2pt);\end{scope}\node at (6.699999999999999,1.35) {接平：仍严格增};\begin{scope}[shift={(11.299999999999999,3.1)},x=0.65cm,y=0.65cm]\draw[->,gray] (-2,0)--(2.3,0) node[right]{$x$};\draw[->,gray] (0,-2)--(0,4.3) node[above]{$y$};\begin{scope}\clip (-2,-2) rectangle (2,4);\draw[blue,thick,domain=-2:0,samples=80] plot(\x,{\x+2});\draw[orange!85!black,thick,domain=0.0001:2,samples=100] plot(\x,{\x});\end{scope}\draw[draw=blue,fill=blue] (0,2) circle(2pt);\draw[draw=orange!85!black,fill=white] (0,0) circle(2pt);\end{scope}\node at (11.299999999999999,1.35) {向下跳：不递增};
\node[text width=12.8cm,align=center] at (7.1,.65) {蓝实心取，橙空心不取；接平时端点重合，按实心取。};
\end{tikzpicture}\end{center}



\textbf{要解决的问题：}两条分别向上走的曲线，拼接后不一定还在向上走。\textbf{不仅要检查每一段，还要检查两段交接的地方。}

设
\[
 f(x)=\begin{cases}g(x),&x\le x_0,\\ h(x),&x>x_0.\end{cases}
\]
这里 $x_0$ 是分界点；$f(x_0)=g(x_0)$，因为等号放在左段。

\textbf{先记最容易懂、又最不容易用错的版本：}
\begin{enumerate}[leftmargin=*,itemsep=0.25em]
\item \textbf{整体严格递增：}左段严格递增；右段严格递增；而且\emph{任意一个}右段的函数值，都必须\emph{大于}左段在分界点的值：
\[
  \boxed{\ g(x_0)<h(x)\quad (\text{对所有 }x>x_0)\ }.
\]
\item \textbf{整体严格递减：}左段严格递减；右段严格递减；而且
\[
  \boxed{\ g(x_0)>h(x)\quad (\text{对所有 }x>x_0)\ }.
\]
\end{enumerate}
上面两个判断，在已知左右两段分别严格递增或递减的情况下，既是必须满足的，也是足够保证整体严格单调的。\textbf{这个版本不用求极限，始终可以直接理解。}

\medskip
\end{statementbox}
\clearpage
\begin{statementbox}
\textbf{高考题最常用的快速写法：}
如果右段 $h(x)$ 在 $x$ 从右边接近 $x_0$ 时，函数值会接近某个\textbf{有限的数} $L$，可记为
\[
 L=\lim_{x\to x_0^+}h(x)
 \quad\text{（就是“右边刚靠近分界点时的高度”）}.
\]
于是
\[
\boxed{\begin{aligned}
\text{整体严格递增}\quad&\Longleftrightarrow\quad
\begin{cases}\text{左段严格递增},\\\text{右段严格递增},\\g(x_0)\le L;\end{cases}\\[5pt]
\text{整体严格递减}\quad&\Longleftrightarrow\quad
\begin{cases}\text{左段严格递减},\\\text{右段严格递减},\\g(x_0)\ge L.\end{cases}
\end{aligned}}
\]

\textbf{特别容易混淆的地方：}要求\emph{严格}递增，交接处却用 $\le$，\textbf{允许等号}；递减时用 $\ge$，也允许等号。这不矛盾，因为右段\emph{不包含}分界点，且右段本身在严格变化。

\textbf{适用范围：}左右两段在各自区间都有定义；若使用 $L$ 的快捷写法，还要求这个 $L$ 是有限数。若求不出有限的 $L$，直接使用上面的“和所有右段函数值比较”版本，不要硬套接口公式。

\textbf{用词约定：}本讲义把“严格递增”理解为：只要 $x_1<x_2$，必有 $f(x_1)<f(x_2)$。若题目只要求“不下降”（允许出现水平的一段），左右两段分别改成“不下降”，交接处的 $\le$ 不变；递减同理。
\end{statementbox}
